\documentclass{article}

\usepackage{fullpage}
\usepackage{url}
\usepackage{graphicx}
\usepackage{amsmath}
\usepackage{physics}
\DeclareMathOperator{\erfi}{erfi}

\title{General Realtivity Excercise 1}
\author{Idan Stark}
\date{\today}

\begin{document}
\maketitle
\section{Warmup}
\subsection{The Metric is symmetric}
First of all, the metric tensor must be a symmetric one. This can be seen from the symmetry of the definition. Assuming $g_{ij} \ne g_{ji}$, we get
\begin{equation}
    ds^2 = g_{ij} dx^i dx^j \ne g_{ji} dx^i dx^j = ds^2
\end{equation}
Which of course can't be the case, hense the metric is always symmetric.
\subsection{A non-uniform metric}
The most simple example of a non-uniform metric is the metric of the 2d Eucledean space in polar coordinates:
\begin{equation}
    g^{ij} = \begin{pmatrix}1 & \\ & r^2\end{pmatrix}
\end{equation}
\subsection{Inveriance of $ds^2$ under change of coordinates}
Let there be two sets of coordinates $x^\mu$ and ${x'}^\mu$, and the metric is defined for the first set of coordinates such that
\begin{equation}
    ds^2 = g_{\mu\nu}dx^{\mu}dx^{\nu}
\end{equation}
The transformation between the two sets of coordinates is given by the Jacobian:
\begin{equation}
\begin{gathered}
    J^\mu_\nu = \frac{\partial x'^\mu}{\partial x^\nu} \qquad
    \left(J^{-1}\right)^\mu_\nu = \frac{\partial x^\mu}{\partial x'^\nu}
    \\
    dx'^\mu = J^\mu_\nu dx^\nu \\
    g'_{\mu\nu} = \left(J^{-1}\right)^\sigma_\mu \left(J^{-1}\right)^\tau_\nu g_{\sigma\tau}
\end{gathered}
\end{equation}
So, using the property of the Jaccobian
\begin{equation}
    J^\mu_\alpha \left(J^{-1}\right)^\beta_\mu = \frac{\partial x'^\mu}{\partial x^\alpha} \frac{\partial x^\beta}{\partial x'^\mu} = 
    \frac{\partial x^\beta}{\partial x^\alpha} = \delta^\beta_\alpha
\end{equation}
The norm in the $x'^\mu$ coordinate system, is
\begin{equation}
    ds'^2 = g'_{\mu\nu}dx'^\mu dx'^\nu = \left(J^{-1}\right)^\sigma_\mu \left(J^{-1}\right)^\tau_\nu g_{\sigma\tau} J_\alpha^\mu dx^\alpha J_\beta^\nu dx^\beta = \delta^\sigma_{\alpha} \delta^\tau_\beta g_{\sigma\tau} dx^\alpha dx^\beta = g_{\sigma\tau} dx^\sigma dx^\tau = ds^2
\end{equation}
\subsection{Another example for the usage of a non-measurable variable}
A good example for that would come from the ladder operators in QM. In addition to the x and p operators, we create another form of looking at the data, which is unobservable but useful in certain situations.
\section{Properties of the Lorentz Transformation}
\subsection{The Lorentz Group}
\subsubsection{Definition}
A Lorentz transformation $\Lambda$ is a linear transformation that preserves the norm $ds^2 = c^2dt^2 - d\vec{x}^2$.
\subsubsection{Multiplication}
Let $\Lambda_\mu^\nu$ and ${\Lambda'}_\mu^\nu$ be two Lorentz transofrmations. Their product is also linear, since the product of the two linear transformation is linear.
\newline
Now, let $v^\mu$ be a vector and $ds^2$ its squared norm. By the definition, the vector ${v'}^\mu = \Lambda_nu^\mu v^\nu$ has the same norm $ds^2$, and again, by the definition, the vector ${v''}^\mu = {\Lambda'}^\mu_\nu \Lambda^\nu_\sigma v^\sigma$ has the same norm as ${v'}^\mu$ and $v^\mu$. This means that the product of the two transformations $\Lambda' \circ \Lambda$ also preserves the norm and therefore a Lorentz transformation.
\subsubsection{The Lorentz Group}
To demonstrate that the set of all Lorentz transformations over a given Minkowski space is a group, we must show that it satisfies the four group axioms.
\newline
The first axiom is closure: the product of two items of the group is also an item of the group. We demonstrated that in section 2.1.2.
\newline
The second axiom is associativity, which is an inherit part of linear transformation product and therefore needs no additional proof.
\newline
The third axiom is the existance of an identity element. It's easy to see that the identity matrix follows the definition of a Lorentz transformation, and therefor is a part of the group.
\newline
The last axiom is the existance of an inverse element. Since all Lorentz transformation are linear, we can represent them in a generic base. Because the transformations preserve the Minkowski norm, their determinant must be $\pm1$ and not zero. This means they must be invertable.
\newline
These four conditions are the four axioms of a gropu, meaning they are enough to proove that the set of all Lorentz Transformation is a group.
\subsubsection{Proper and Improper Lorentz Transformations}
A Lorentz Transformation is proper if $\det(\Lambda) = 1$ and improper if $\det(\Lambda) = -1$. The product of two proper transformation is also proper, since determinant of a product of matrices is the product of the determinants:
\begin{equation}
    \det(\Lambda_{p1}\circ\Lambda_{p2}) = \det(\Lambda_{p1})\det(\Lambda_{p2}) = 1\cdot1 = 1
\end{equation}
Which means that the set of all proper Lorentz transformation over a given Minkowski space is closed under multiplication, and thus is a subset of the Lorentz group.
\newline
An example of an improper Lorentz Transformation would be:
\begin{equation}
    \Lambda = \begin{pmatrix}
        -1 & & & \\
        & 1 & & \\
        & & 1 & \\
        & & & 1
    \end{pmatrix}
\end{equation}
This is the time reversal Transformation, and it's a Lorentz Transformation since for an infinitisimal vecotr $dx^\mu = (dt, d\vec{x})$ the transformed vector is ${x'}^\mu = \Lambda^\mu_\nu x^\nu = (-dt, d\vec{x})$, and the new norm is ${ds'}^2 = (-dt)^2 + d\vec{x}^2 = dt^2 + d\vec{x}^2 = ds^2$.
\subsubsection{Orthochronous and Antichronous Lorentz Transformation}
An Orthochronous Lorentz Transformation is one such that
\begin{equation}
    \Lambda(e_0)\cdot e_0 > 0
\end{equation}
Where $e_0$ is a future directed timelike unit vector. Using the definition of the inner product and writing in index notation we get
\begin{equation}
    \eta_{\alpha\beta}\Lambda^\alpha_\gamma e^\gamma_0e_0^\beta > 0
\end{equation}
\subsection{On Boosts}
\subsubsection{Boost in an arbitrary direction}
To get a boost in an arbitrary direction, we simply use a rotation matrix with angles $\theta$ and $\varphi$ turning the system around the $x^2$ and $x^3$ axes respectivly:
\begin{equation}
    R_y(\theta) = \begin{pmatrix}
        1 &  &  & \\
          & \cos\theta & &  -\sin\theta \\
          &  & 1 & \\
          & \sin\theta &  &  \cos\theta
    \end{pmatrix}
    \qquad
    R_z(\varphi) = \begin{pmatrix}
        1 & & & \\
        & \cos\varphi & -\sin\varphi & \\
        & \sin\varphi & \cos\varphi & \\
        & & & 1
    \end{pmatrix}
\end{equation}
Using these matrices to rotate the $\vec{\beta} = \beta\hat{x}$ vector and the $\Lambda$ Lorentz transformation along the x axis, around the y axis and then the z axis, we get:
\begin{equation}
\begin{gathered}
    \beta' =  R_3(\varphi) R_2(\theta) \vec{\beta} = \beta\begin{pmatrix}
        \\
        \cos \theta \cos \varphi \\
        \sin \varphi \\
        \sin \theta \cos \varphi \\
    \end{pmatrix}
    \equiv
    \begin{pmatrix}
        \\
        \beta^1 \\
        \beta^2 \\
        \beta^3 \\
    \end{pmatrix}
    \\
    \Lambda' = R_2(\theta) R_3(\varphi) \Lambda R_3(-\varphi) R_2(-\theta) = 
    \\
    = \begin{pmatrix}
            \gamma &
            -\gamma\beta\cos\theta\cos\varphi &
            -\gamma\beta\sin\varphi &
            -\gamma\beta\sin\theta\cos\varphi
        \\
            -\gamma\beta\cos\theta\cos\varphi &
            1 + (\gamma - 1)\cos^2\theta\cos^2\varphi &
            (\gamma - 1)\cos\theta\cos\varphi\sin\varphi &
            (\gamma - 1)\cos^2\varphi\cos\theta\sin\theta
        \\
            -\gamma\beta\sin\varphi &
            (\gamma - 1)\cos\theta\cos\varphi\sin\varphi &
            1 + (\gamma - 1)\sin^2\varphi &
            (\gamma - 1)\sin\theta\cos\varphi\sin\varphi
        \\
            -\gamma\beta\sin\theta\cos\varphi &
            (\gamma\ - 1)\cos^2\varphi\cos\theta\sin\theta &
            (\gamma - 1) \sin\theta\cos\varphi\sin\varphi &
            1 + (\gamma - 1)\sin^2\theta\cos^2\varphi
    \end{pmatrix}
\end{gathered}
\end{equation}
We can replace the angles in the matrix with the respected values of $\beta$.
\begin{equation}
    \Lambda' = \begin{pmatrix}
            \gamma &
            -\gamma\beta^1 &
            -\gamma\beta^2 &
            -\gamma\beta^3
        \\
            -\gamma\beta^1 &
            1 + (\gamma - 1)\left(\frac{\beta^1}{\beta}\right)^2 &
            (\gamma - 1)\frac{\beta^1\beta^2}{(\beta)^2} &
            (\gamma - 1)\frac{\beta^1\beta^3}{(\beta)^2}
        \\
            -\gamma\beta^2 &
            (\gamma - 1)\frac{\beta^1\beta^2}{(\beta)^2} &
            1 + (\gamma - 1)\left(\frac{\beta^2}{\beta}\right)^2 &
            (\gamma - 1)\frac{\beta^2\beta^3}{(\beta)^2}
        \\
            -\gamma\beta^3 &
            (\gamma\ - 1)\frac{\beta^1\beta^3}{(\beta)^2} &
            (\gamma - 1) \frac{\beta^2\beta^3}{(\beta)^2} &
            1 + (\gamma - 1)\left(\frac{\beta^3}{\beta}\right)^2
    \end{pmatrix}
\end{equation}
And writing the components gives us:
\begin{equation}
\begin{gathered}
    \hat{x}^0 = \gamma(x^0 - \vec{\beta}\cdot\vec{x})
    \\
    \hat{\vec{x}} = \vec{x} + \frac{\gamma - 1}{\beta^2}\left(\vec{\beta}\cdot\vec{x}\right)\vec{\beta} - \gamma\vec{\beta}x^0
\end{gathered}
\end{equation}
\subsubsection{Relativistic velocity addition of velocities}
Our new velocity will be $\hat{\vec{u}} = d\hat{\vec{x}}/d\hat{t}$ where $\hat{t} = \hat{x}^0/c$. Using the last section, we can write the infinitisimal length in each coordinate
\begin{equation}
\begin{gathered}
    d\hat{x}^0 = \gamma dx^0 - \vec{\beta}\cdot d\vec{x} = \gamma\left(1 - \vec{\beta}\cdot \frac{\vec{u}}{c}\right)dx^0
    \\
    d\hat{\vec{x}} = d\vec{x} + \frac{\gamma - 1}{\beta^2}\left(\vec{\beta}\cdot d\vec{x}\right)\vec{\beta} - \gamma\vec{\beta}dx^0 = \frac{1}{c}\left(\vec{u} + \frac{\gamma - 1}{\beta^2}\left(\vec{\beta}\cdot \vec{u}\right)\vec{\beta} - \gamma c\vec{\beta}\right)dx^0
\end{gathered}
\end{equation}
Splitting $\vec{u}$ into $\vec{u} = \vec{u}_{\|} + \vec{u}_\perp$, we get
\begin{equation}
\begin{gathered}
    d\hat{x}^0 = \gamma\left(1 - \beta \frac{u_{\|}}{c}\right)dx^0
    \\
    d\hat{\vec{x}} = \frac{1}{c}\left(\vec{u}_{\perp} + \gamma\vec{u}_{\|} - \gamma c\vec{\beta}\right)dx^0
\end{gathered}
\end{equation}
And the velocity is
\begin{equation}
    \hat{\vec{u}} = c\frac{d\hat{\vec{x}}}{d\hat{x}^0} = \frac{\vec{u}_{\perp} + \gamma\vec{u}_{\|} - \gamma c\vec{\beta}}{\gamma\left(1 - \beta \frac{u_{\|}}{c}\right)}
\end{equation}
And for $c\rightarrow\infty$, which gives $\beta = 0$, $\vec{\beta} c = \vec{v}$, and $\gamma = 1$, the denominator becomes 1 and we get:
\begin{equation}
    \hat{\vec{u}} = \vec{u}_{\perp} + \vec{u}_{\|} - \vec{v} = \vec{u} - \vec{v}
\end{equation}
\section{Rocket}
\subsection{Non-Relativistic Rocket}
\subsubsection{Instantaneously Infinitesimal Change in the Rocket’s Velocity in its reference frame}
Because the burning rate is constant, the masses of the rocket and the gasses expulted as a function of time is
\begin{equation}
\begin{gathered}
    M_r(t) = M_i + \frac{M_f - M_i}{T}t \\
    M_g(t) = \frac{M_i - M_f}{T}t
\end{gathered}
\end{equation}
Since we assume an inertial frame of reference, we can use the preservation of momentum. The infinitisimal momentum leaving the rocket is
\begin{equation}
    dp = dM_g(t) v_e + M_g(t) dv_e = \frac{M_i - M_f}{T} v_e dt
\end{equation}
Because of the preservation of momentum, this is also the infinitisimal increase in the rocket's momentum to the opposite direction:
\begin{equation}
    dp = dM_r(t) v_r + M_r(t) dv_r = \frac{M_f - M_i}{T} v_r dt + \left(M_i + \frac{M_f - M_i}{T}t\right) dv_r = \frac{M_i - M_f}{T} v_e dt
\end{equation}
Assigning $v_r = 0$, since the rocket is at rest, and solving for $dv_r$, we get:
\begin{equation}
    dv_r = \frac{M_i - M_f}{M_i T + (M_f - M_i)t} v_e dt
\end{equation}
\subsubsection{The Acceleration $a(t)$}
Simply dividing equation 22 by $dt$ gives
\begin{equation}
    a(t) = \frac{M_i - M_f}{M_i T + (M_f - M_i)t} v_e = \frac{\mu}{1 - \mu\tau} \frac{v_e}{T}
\end{equation}
Where
\begin{equation}
    \mu = \frac{M_i - M_f}{M_i} \qquad \tau = \frac{t}{T}
\end{equation}
\subsubsection{Final Velocity}
By integrating over equation 22 from $t = 0$ to $t = T$ on the right and from $v = 0$ to $v = v_f$ on the left, we get
\begin{equation}
    v_f = v_e \int_{0}^{1} \frac{\mu}{1 - \mu\tau} d\tau = - v_e \left[\ln(1 - \mu\tau) \right]_0^1 = v_e \ln\frac{1}{1-\mu} = v_e ln\frac{M_i}{M_f}
\end{equation}
It can be seen that since the rocket's mass continues to decrease, the constant expulsion of gas, while increasing its momentum by a constant amount, increases the rocket's velocity by a uniformally increasing amount, causing it to eventually surpass the gas expulsion velocity.
\subsection{Relativistic Rocket}
\subsubsection{Moving between reference frames}
When moving between inertial reference frames, the acceleration follows the Lorentz transformation, but also follows the formula in equation 23. For instance, moving to the inertial refrence frame where $v_e = 0$ would cause the acceleration to be zero as well, which of course cannot be the case.
\subsubsection{Momentum and Energy Conservation in a generic reference frame}
In a generic inertial reference frame, where the velocities of the rocket and the gas are $v_1$ and $v_2$ respectivly, and their masses are $m_1$ and $dm_2$ respectivly, the conservation of energy and momentum can be formulated as:
\begin{equation}
\begin{gathered}
    m_1 v_1 = (m_1 + dm_1) (v_1 + dv_1) + dm_2 v_2 \\
    \frac{1}{2} m_1 v_1^2 = \frac{1}{2}(m_1 + dm_1) (v_1 + dv_1)^2 + \frac{1}{2}dm_2 v_2^2
\end{gathered}
\end{equation}
Or
\begin{equation}
\begin{gathered}
    0 = m_1 dv_1 + dm_1 v_1 + dm_1 dv_1 + dm_2 v_2 \\
    0 = 2 m_1 v_1 dv_1 + m_1 dv_1^2 + dm_1 v_1^2 + 2 dm_1 v_1 dv_1 + dm_1 dv_1^2 + dm_2 v_2^2 
\end{gathered}
\end{equation} 
\subsubsection{The masses $dm_1$ and $dm_2$ as a function of $m_1$, $v_1$, $v_2$, and $dv_1$}
Solving equation 27 we get
\begin{equation}
\begin{gathered}
    (v_1 + dv_1) dm_1 + v_2 dm_2 = -m_1 dv_1 \\
    (v_1 + dv_1)^2dm_1 + v_2^2 dm_2 = -2m_1v_1dv_1 + m_1dv_1^2
    \\
    \\
    dm_2 = \frac{2dv_1 - v_1}{v_2(v_2 - v_1 - dv_1)}m_1 dv_1 \qquad
    dm_1 = -\frac{v_2 - 2v_1 + dv_1}{(v_1 + dv_1)(v_2 - v_1 - dv_1)}m_1 dv_1
\end{gathered}
\end{equation}
\end{document}